\title{A Large-scale Study of Representation Learning with the Task Adaptation Benchmark}

\begin{abstract}
Representation learning promises to unlock deep learning for the long tail of vision tasks without expensive labelled datasets. Yet, the absence of a unified evaluation for general visual representations hinders progress. Popular protocols are often too constrained (linear classification), limited in diversity (ImageNet, CIFAR, Pascal-VOC), or only weakly related to representation quality (ELBO, reconstruction error). We present the Visual Task Adaptation Benchmark (VTAB), which defines good representations as those that adapt to \emph{diverse}, \emph{unseen} tasks with \emph{few examples}. With VTAB, we conduct a large-scale study of many popular publicly-available representation learning algorithms. We carefully control confounders such as architecture and tuning budget. We address questions like: How effective are ImageNet representations beyond standard natural datasets? How do representations trained via generative and discriminative models compare? To what extent can self-supervision replace labels? And, how close are we to general visual representations?
\end{abstract}

\section{Introduction}

Distributed representations learned from raw pixels have enabled unprecedented performance on many visual understanding tasks. Hand-crafted features have been replaced with hand-annotated datasets, with thousands to millions of examples~\citep{cifar10,imagenet}. By contrast, humans learn a wide range vision tasks using just a few examples per task. A key research challenge is to close this gap in sample efficiency, and unlock deep learning for the long tail of problems without many labels.

Improving sample efficiency has been approached from many angles: few-shot learning~\citep{li2006}, transfer learning~\citep{pan2009survey}, domain adaptation~\citep{wang2018deep}, and representation learning~\citep{bengio2013representation}. Representation learning is studied in many contexts: supervised pre-training~\citep{razavian2014}, self-supervised learning~\citep{doersch2015unsupervised}, semi-supervised learning~\citep{chapelle2009semi}, generative modeling~\citep{donahue2016adversarial}, and disentanglement learning~\citep{higgins2017beta}. Each sub-domain has its own evaluation protocol, and the lack of a common benchmark impedes progress. Benchmarks have been critical in other sub-fields, such as RL~\citep{atari}, image classification~\citep{imagenet}, and NLP~\citep{glue}. Inspired by these successes, we propose a benchmark with similar principles: (i) minimal constraints to encourage creativity, (ii) a focus on practical considerations, and (iii) make it challenging.

We present the Visual Task Adaptation Benchmark (VTAB). Using VTAB, we perform the first extensive cross sub-field study of representation learning. VTAB defines a good representation as one that can be used to solve many diverse \emph{previously unseen} tasks with the \emph{fewest possible labels}. We consider low sample complexity to be the key objective of representation learning. Task diversity is also crucial to assess generality. Therefore, VTAB goes beyond standard natural tasks, and includes those related to sensorimotor control, medical imaging, and scene understanding.

In our study, we investigate many representation learning algorithms, pre-trained on ImageNet. We carefully control confounding factors such as tuning budget, architecture, and pre-training data. This study quantifies existing intuitions and reveals new insights: (i)~Supervised ImageNet pre-training yields excellent representations for natural image classification tasks. Interestingly, it also yields a smaller, but consistent improvement on specialized tasks (e.g. medical imaging). However, these representations are extremely limited for tasks that require structured understanding. (ii)~Self-supervised is less effective than supervised learning overall, but surprisingly, can improve structured understanding. (iii)~Combining supervision and self-supervision is effective, and to a large extent, self-supervision can replace, or compliment labels. (iv)~Discriminative representations appear more effective than those trained as part of a generative model, with the exception of adversarially trained encoders. (v)~GANs perform relatively better on data similar to their pre-training source (here, ImageNet), but worse on other tasks. (vi)~Evaluation using a linear classifier leads to poorer transfer and different conclusions.

\section{The Visual Task Adaptation Benchmark}

We seek algorithms that perform well on a wide variety of unseen visual understanding tasks with few labels per task.\footnote{
In some practical settings additional \emph{unlabelled} data may be available for unseen tasks in addition to the labelled data. We omit this setting, leaving it to future work.} We first formalize this objective and then specify a practical benchmarking procedure to measure progress.

A \emph{dataset} $D^n$ is a set of $n$ instances $\{({\bm{x}}_i,\, {\bm{y}}_i)\}_{i=1}^n$ with observations  ${\bm{x}}_i \in X$ and labels ${\bm{y}}_i \in Y$. A \emph{prediction function} is any mapping $F: X \to Y$ (e.g.\ a classifier). A (learning) \emph{algorithm}, $\mathcal{A}$, takes as input a dataset and outputs a prediction function. For example, $\mathcal{A}$ may be a pre-trained network coupled with a training mechanism. An \emph{evaluation procedure}, $\mathcal{E}_T$, takes $F$ and outputs a scalar measuring $F$'s performance (e.g.\ test-set accuracy). We seek the algorithm that maximizes the expected performance over a distribution of tasks $P_T$, where a task $T$ is a tuple containing a task-specific dataset distribution $D_T$ and evaluation procedure. Given only $n$ samples per task we want to maximize:

\begin{align}
\textsc{score}_n(\mathcal{A}) = \mathbb{E}_{T \sim P_T}\; \mathcal{E}_T\left[ \mathcal{A}(D_T^n) \right],\label{eq:ta}
\end{align}

This general formulation requires a few clarifications. First, the task distribution needs to be appropriate; we desire a spectrum that covers tasks solvable by a vision algorithm with human-like capabilities. Second, $n$ may be varied to measure an algorithm's sample complexity. In practice we choose $n$ to be similar to a modest labelling budget (\ref{sec:setup}).
Third, we assume that $P_T$ is known, and we aim to build an algorithm with the best inductive biases for solving samples from it. We desire an ``open world'', where evaluation data is extremely diverse and can always be freshly sampled. Unfortunately, in practice the benchmark must contain a finite number of test tasks. Therefore, we must ensure that the algorithms are not pre-exposed to specific evaluation samples, as described below.

\subsection{A Practical Benchmark} We now describe the Visual Task Adaptation Benchmark (VTAB), which is designed to be the best possible proxy for \ref{eq:ta}.
\ref{fig:vtab_protocol} present an overview.

\textbf{Task Distribution}\quad
We aim for universal visual understanding, so we informally define $P_T$ as ``Tasks that a human can solve, from visual input alone.''. We validate this empirically, see \ref{app:human}.
Intuitively, such tasks should benefit from visual representations learned by observing and interacting with the natural world. The second clause eliminates tasks that require external knowledge that is (currently) unreasonable for a vision algorithm to acquire -- an extreme example would be to classify objects grouped by their spelling in a natural language.
Section~\ref{sec:tasks} details the samples.

\textbf{Expectation Over Tasks}\quad
We approximate the expectation over tasks by an empirical average over a number of hand-picked samples. Ideally, we would sample a new task for each evaluation, as is possible in procedural environments, e.g.~\citep{finn2017}. However, real-world vision tasks are expensive to collect, so we define a fixed, representative set of samples from $P_{T}$. While fixing the set reduces variance, it introduces a risk of meta-overfitting.

\textbf{Mitigating Meta-Overfitting}\quad
To reduce meta-overfitting, we treat the evaluation tasks like a test set, and consider them unseen. Therefore, algorithms that use pre-training must not pre-train on any of the evaluation tasks (even their unlabelled images). Upstream training should provide useful inductive biases for any draw from $P_{T}$, so should not use test samples. Fortunately, despite numerous test re-evaluations on popular ML benchmarks, such as ImageNet, progress seems to transfer to new data~\citep{recht18cifar,recht19imagenet,kornblith2018better}.

\textbf{Unified Implementation}\quad
With many diverse tasks, usage could become impractical. As discussed, the algorithms must have no prior knowledge of the downstream tasks; while it is permitted to run a hyperparameter search on each task, the search space cannot be task-dependent. To get meaningful results, we need to define hyperparameter searches that work well across the benchmark.

For this, we convert all tasks into classification problems. For example, a detection task requiring localization of an object can be mapped to classification of the $(x,y,z)$ coordinates. With a homogeneous task interface, we may control for possible confounding factors. For example, we may use the same architecture and hyperparameter sweep everywhere. Not all tasks can be efficiently modeled as image-level classification, such as those requiring per-pixel predictions. Nonetheless, we design the tasks such that success on VTAB requires learning of diverse set of visual features: object identification, scene classification, pathology detection, counting, localization, and 3D geometry.

\subsection{Tasks\label{sec:tasks}}
VTAB contains 19  tasks which cover a broad spectrum of domains and semantics. \ref{app:tasks} contains details. These are grouped into three sets: \taskNatural{}, \taskSpecialized{}, and \taskStructured{}.

The \taskNatural{} group represents classical vision problems. These tasks contain natural images captured using standard cameras. The classes may represent generic, fine-grained, or abstract objects. The group includes: Caltech101, CIFAR-100, DTD, Flowers102, Pets, Sun397, and SVHN.

The \taskSpecialized{} group also contains images of the world, but captured through specialist equipment. These images have different invariances to those in the \taskNatural{} tasks. Nonetheless, humans recognize the structures therein, thus generic visual representations should also capture the visual concepts. We have two sub-groups: remote sensing, and medical. Remote sensing includes Resisc45 and EuroSAT: aerial images of the earth captured using satellites or aerial photography. Medical includes Patch Camelyon, metastases detection from microscopy images, and Diabetic Retinopathy, retinopathy classification from fundus images.

The \taskStructured{} group assesses comprehension of the structure of a scene, for example, object counting, or 3D depth prediction. Most of these tasks are generated from simulated environments, whose structure is easy for a human to determine, but whose domain differs greatly to datasets like ImageNet. These tasks are intended as a step towards useful representations for perceptual control. We include:
\emph{Clevr}: Simple shapes rendered in a 3D scene, with two tasks: counting and depth prediction.
\emph{dSprites}: Simple black/white shapes rendered in 2D, with two tasks: location and orientation prediction.
\emph{SmallNORB}: Artificial objects viewed under varying conditions, with two tasks: object-azimuth and camera-elevation prediction.
\emph{DMLab}: Frames from a rendered 3D maze. The task involves predicting the time for a pre-trained RL agent to navigate to an object.
\emph{KITTI}: frames captured from a car driver's perspective. We predict the depth of the nearest vehicle.

\subsection{Representation and Transfer Learning}

Success on VTAB, and ultimately optimizing~\ref{eq:ta}, requires some knowledge of $P_{T}$. For example, CNN architectures have a useful inductive bias for vision. While these biases are usually manually designed (e.g. through architecture choice), they can also be learned. For example by learning: architectures~\citet{zoph2017}, optimization algorithms~\citet{bello2017neural}, initialization distributions~\citet{raghu2019},
or pre-trained data representations.

Whilst many strategies merit investigation, we focus on representation learning. Human visual perception is refined through years of observation and interaction with the world, resulting in a system that solves new tasks with few in-domain labels. Likewise, we aim to pre-train a network that extracts useful features, or representations, from raw data.

\textbf{Transfer Strategy}\quad
The representations are not pre-trained on the evaluation tasks themselves (which VTAB forbids), so they must be adapted to solve the new tasks (using limited in-domain data). The simplest adaptation strategy in deep representation learning is to first pre-train the network, then \emph{freeze} the network's weights and train another -- usually smaller -- model on top. When the upstream and downstream datasets differ significantly,
\emph{fine-tuning} the original weights is more effective~\citep{yosinski2014,kornblith2018better}. VTAB does not constrain the transfer strategy; here we use fine-tuning as it tends to perform best.

\textbf{Upstream Training}\quad
Representation learning literature often focuses on unsupervised learning which may be applied to any dataset. However, supervised data, where available, can yield good representations. Indeed, the most popular models used in practice are pre-trained on ImageNet labels~\citep[and refs therein]{huh2016makes}. VTAB \emph{does not constrain the type of data} used for pre-training.

\section{Large-Scale Study\label{sec:experiments}}

With VTAB, we evaluate many popular, publicly-available representation learning algorithms across the different types of tasks. Here, we study the pre-training losses, and therefore, control for other factors such as architecture, preprocessing, transfer hyperparameters, and pre-training data. Of course, VTAB may be used subsequently to study and improve these aspects as well. Finally, we provide additional analysis of various of VTAB's design choices and assess other evaluation protocols.

\subsection{Setup\label{sec:setup}}

\textbf{Downstream Data Size}\quad
Section~\ref{sec:tasks} describes the tasks. VTAB aims to assess adaptation with limited data, so we evaluate primarily using \num{1000} labelled examples per task (called VTAB-1k). We define train (\num{800} samples) and validation (\num{200} samples) sets for users' convenience. We emphasize that to avoid meta-overfitting, one should not use extra images for hyperparameter selection, but may use the \num{1000} examples in any manner. We also use the full datasets (VTAB-full). This allows us to assess the value of representation learning as in-domain data increases, and to check performance against prior art.

\textbf{Representation Learning Algorithms}\quad
We evaluate supervised, semi-supervised, self-supervised, and generative models, selecting popular, public methods from each class. We also compare to ``from-scratch'' models, which use no pre-training, yielding a total of 18 algorithms. All models (except ``from-scratch'') are pre-trained on ImageNet (1.28M labelled images). Some use the ImageNet labels, others use some or none of the labels.

We train two supervised models: one using $100\%$ of the ImageNet labels (\textsc{Sup-100\%}), and one using $10\%$ (\textsc{Sup-10\%}).
For self-supervised learning, we include both image-based and patch-based models. The image-based models include \textsc{rotation}~\citep{gidaris2018unsupervised} and \textsc{Exemplar}~\citep{dosovitskiy2014exemplar}.
Patch-based include
\textsc{Relative Patch Location}~\citep{doersch2015unsupervised} and \textsc{Jigsaw}~\citep{noroozi2016unsupervised}.
We use the public implementations of~\citet{kolesnikov2019revisiting}.
The patch-based models are converted into image-level classifiers by averaging the representations from a $3\times3$ grid of patches (see \ref{app:architectures}).
The semi-supervised models are trained using 10\% of the ImageNet labels with an auxiliary loss on all of the data, see~\citep{zhai2019s4l}. We use either rotation (\textsc{Sup-Rotation-10\%}) or Exemplar (\textsc{Sup-Exemplar-10\%}) auxiliary losses. These models can also use all of the labels, denoted \textsc{Sup-Rotation-100\%} and \textsc{Sup-Exemplar-100\%}.
For the generative models, we evaluate GANs and VAEs. As is common, we take the GAN's image representations from the discriminator, replacing the final linear layer with a classification layer. We use both the label-conditional and unconditional BigGAN discriminators~\citep{brock2018large} (\textsc{Cond-BigGAN} and \textsc{Uncond-BigGAN}) using the implementations of~\citet{chen2019self,lucic2019high}.
We also evaluate the encoder from \textsc{BigBiGAN}~\citep{donahue2019large}.
For the autoencoders, we use the encoder as the representation. We evaluate VAEs~\citep{kingma2013auto}, and WAEs with three distribution matching losses: GAN, MMD~\citep{tolstikhin2017wasserstein}, and UKL~\citep{rubenstein2019practical}.

\textbf{Transfer Algorithm}\quad
VTAB permits any method of transfer. Here, we fine-tuning the entire network on the task data, since that performs best~\citep{kornblith2018better}. It is also popular to add a linear model to a frozen network, so we compare to this approach in \ref{sec:linear}.

\textbf{Hyperparameters}\quad
Upstream, we contol the data (ImageNet) and architecture. We find that bigger architectures perform better on VTAB (\ref{app:scale-architecture} and~\citet{kolesnikov2019revisiting}).
For comparability, we use ResNet50-v2, or similar, in this study. For the supervised and semi-supervised methods we use the architecture in~\citet{he2016identity}.
For GANs and auto-encoders we use very similar deep ResNets, but with appropriate modifications to pre-train them successfully as a generative model.
\ref{app:architectures} contains details.

Downstream, the transfer hyperparameters influence performance, and different datasets require different settings. We run VTAB in two modes: \emph{lightweight} and \emph{heavyweight}.
Lightweight mode performs a restricted per-task hyperparameter search. Many hyperparameters are fixed, including optimizer, batch size, pre-processing, and weight decay. For each task, the lightweight mode sweeps: 2 initial learning rates, and 2 learning rate schedules. See \ref{app:hyperparameter_sweep} for details. We choose short schedules to limit cost, but show in \ref{sec:analysis} that these yield near-optimal performance.

In heavyweight mode we perform a large random search over learning rate, schedule, optimizers, batch size, train pre-processing functions, evaluation pre-processing, and weight decay. We include longer training schedules and higher resolution images. This mode is used to understand the impact of a larger computational budget on the performance, and establish performance upper bounds on the methods used.
\ref{app:hyperparameter_sweep} contains details.

We perform our main study in lightweight mode. In \ref{sec:heavy} we show that although extensive tuning improves all methods, the relative performances are mostly unaffected. For future study using VTAB, we recommend defining a similar lightweight setup (or using the one here) that facilitates fair comparison without high cost. If ample computational resources are available, a heavyweight mode may be used to improve the state-of-the-art.

\textbf{Tuning and Evaluation Protocol}\quad
We perform adaptation on the training set and perform model selection for each task on the pre-defined validation sets. We re-train the best model on the union of training and validation sets and evaluate on the test set. Note, for VTAB-1k we define custom train set with 800 examples and validation set with 200 examples. We run the final evaluation on the test set with three random seeds and report the median score.

\textbf{Metrics}\quad
We evaluate with top-1 accuracy. We consider other metrics in \ref{app:metrics}, and find the conclusions are the same. To aggregate scores across tasks, we take the mean accuracy. We investigate more complex aggregation strategies in Section~\ref{sec:analysis}, but the relative performances are unaffected, so we use mean accuracy for simplicity.

\subsection{VTAB Results}

We run all 18 methods on 19 tasks for both VTAB-1k and VTAB-full using the lightweight sweep.
\ref{fig:all-methods-grouped} shows aggregate performance of each method group (generative, self-sup, etc.) on each VTAB task group (\taskNatural{}, etc.).
\ref{fig:all-methods} presents a breakdown of the performance of each algorithm on each taks group.
\ref{fig:atari} shows two detailed per-task comparisons:
\textsc{Sup-100\%} (ImageNet) versus \textsc{From-Scratch}, and
\textsc{Sup-Rotation-100\%} (best overall on VTAB-1k) versus \textsc{Sup-100\%}.
\ref{app:lightweight},~\ref{tab:lightweight} contains the complete results table.

\textbf{Generative Models}\quad
Overall, these perform worst, \ref{fig:all-methods-grouped}.
The autoencoders perform more poorly than training from scratch. Despite SOTA generative quality on ImageNet, the GANs' discriminators do not provide useful representations. The GAN discriminators appear unstable, with large error bars in \ref{fig:all-methods}.
The \textsc{BigBiGAN} encoder stands out, performing much better -- similar to the best self-supervised models, \ref{fig:all-methods}.
Interestingly, \textsc{BigBiGAN} performs relatively well on natural tasks, but less so on structured tasks.
\textsc{Uncond-GAN} shows a similar pattern. This indicates that GANs fit more strongly to ImageNet's domain (natural images), than self-supervised alternatives. Despite huge advances in generative quality, these models do not yield great representation quality.
Similarly,~\citet{ravuri2019classification}~show that models with high sample quality fail to generate data from which an accurate classifier can be trained. However, the BigBiGAN result holds promise for adversarially-trained encoders, and consideration of downstream representation quality may lead to further progress.

\textbf{Self-supervised}\quad
All self-supervised representations outperform from-scratch training. The best, \textsc{Rotation}, attains $59.6\%$ on VTAB-1k, while \textsc{From-Scratch} attains $42.1\%$ (\ref{app:lightweight}, \ref{tab:lightweight}).
Methods applied to the entire image (\textsc{Rotation}, \textsc{Exemplar}) outperform patch-based methods (\textsc{Jigsaw}, \textsc{Rel.Pat.Loc.}).
However, the patch based methods perform slightly better on DTD (Describable Textures) and Retinopathy (Fundus images), see \ref{app:lightweight}, \ref{tab:lightweight}.
Intuitively, these tasks require sensitivity to local textures, which are indeed captured by patch-based methods. On \taskNatural{} tasks, self-supervision is far behind supervised methods (\ref{fig:all-methods-grouped}).
On \taskSpecialized{}, the performances are similar, and most interestingly, on \taskStructured{} self-supervised methods performs slightly \emph{better} than supervised. This indicates supervised ImageNet models are invariant to useful features required for structured understanding, but self-supervised methods can capture these to some degree.

\textbf{(Semi-)Supervised}\quad
Overall, supervised models perform best. The benefits are most pronounced on the \taskNatural{} tasks, whose domain and semantics are arguably most similar ImageNet classification. However, with self-supervision, the benefit can be attained with fewer labels.
\textsc{Sup-10\%} (120k labels) attains $61.6\%$ on VTAB-1k , while \textsc{Sup-100\%} (1.2M labels) attains $65.6\%$.
With self-supervision, \textsc{Sup-Rotation-10\%} closes $80\%$ of the gap, attaining $64.8\%$.
Recent work has shown in-domain benefits of semi-supervised learning~\citep{henaff2019data,berthelot2019remixmatch}; our study reveals that semi-supervised methods also learn good representations for transfer to new tasks. Finally, additional self-supervision even improves on top of $100\%$ labelled ImageNet, particularly on \taskStructured{} tasks (\ref{fig:atari}.
\textsc{Sup-100\%} attains $65.6\%$, whereas \textsc{Sup-Rotation-100\%} attains $67.5\%$.

\subsection{Heavyweight Hyperparameter Sweeps \label{sec:heavy}}

We evaluate \textsc{From-Scratch} and the best models: \textsc{Sup-100\%}, \textsc{Sup-Rotation-100\%}, and \textsc{Sup-Exemplar-100\%} using the heavyweight hyperparameter search:~\ref{tab:heavyweight_full_test_selection}.
\ref{app:hyperparameter_sweep} shows the selected hyperparameters. As expected, all methods improve significantly. Prior work~\citep{he2018rethinking} shows that with sufficient data and training time, from-scratch training is competitive for detection. However, we observe that across all task groups, pre-trained representations are better than a tuned from-scratch model. On a couple of tasks \textsc{From-Scratch} performs competitively --- Clevr-Dist and sNORB-Elev, which require localization and camera elevation respectively --- indicating that even the best ImageNet representations fail to capture these aspects. The ranking of the best methods is similar, but not identical, with a combination supervision and self-supervision (\textsc{Sup-Exemplar-100\%}) getting the best performance.

We check our results against those recently reported in the literature (\ref{app:heavyweight}, \ref{tab:literature}).
Our results are comparable; behind on highly popular tasks on which complex architectures have been optimized (e.g. CIFAR), but ahead in others. In \ref{app:scale-architecture} we show that simply increasing the architecture size improves VTAB performance.

\subsection{Frozen Features Extractors\label{sec:linear}}

Representations are often evaluated by training a linear layer on a frozen model~\citep{kolesnikov2019revisiting}. ImageNet is often used for linear evaluation, but this is meaningless for the models we consider because many use ImageNet labels for pre-training. However, linear evaluation is also used in a transfer setting~\citep{goyal2019scaling}, so we apply this protocol to the VTAB tasks, and contrast it to fine-tuning. The full protocol resembles the lightweight fine-tuning sweep, see \ref{app:linear}.

\ref{fig:finetune_vs_linear} shows the per-task correlation between linear and fine-tuning.
\ref{app:linear}, \ref{tab:lightweight_linear_test} contains the full table of results. We first note that linear evaluation significantly lowers performance, even when downstream data is limited to $1000$ examples.
\textsc{Sup-100\%} attains $65.6\%$ with fine-tuning (lightweight sweep), but $57.3\%$ with linear. Linear transfer would not by used in practice unless infrastructural constraints required it. Second, \ref{fig:finetune_vs_linear} shows that on many datasets, particularly \taskSpecialized{} and \taskStructured{}, the correlation is low. Linear evaluation may lead to different conclusions, for example, \textsc{Cond-BigGAN} attains $43.3\%$ (1000-examples) on linear, outperforming both \textsc{Rel.Pat.Loc} and \textsc{Jigsaw}.
Yet when fine-tuned these self-supervised methods significantly outperform the GAN discriminator. Another discrepancy is between semi-supervised and supervised.
\textsc{Semi-Rotation-10\%} and \textsc{Semi-Exemplar-10\%} are $1-2\%$ behind \textsc{Sup-100\%} with fine-tuning, but $4-5\%$ behind with linear evaluation. These self-supervised methods extract useful representations, just without linear separability. Previous works~\citep{kornblith2018better,kolesnikov2019revisiting} claim that linear evaluation results are sensitive to additional factors that we do not vary, such as ResNet version or pre-training regularization parameters. Overall linear evaluation is a poor proxy for overall reduced sample complexity.

\subsection{Analysis\label{sec:analysis}}

\textbf{Metrics}\quad
We explore alternatives to using top-1 accuracy, and aggregation of scores accross tasks using the mean. Instead of top-1 accuracy, we used mean-per-class accuracy and Cohen's quadratic kappa. These metrics create only minor ranking differences, leaving the overall picture unchanged. Kendall's ranking correlation with respect to top-1 accuracy is always $>0.97$. Details are in~\ref{app:metrics}.
For aggregation across tasks, we consider seven alternative strategies: (i)~Macro-average across datasets (grouping tasks with the same input data). (ii)~Macro-average across groups (merging tasks in the same group). (iii)~Geometric mean. (iv)~Average rank. (v)~Average rank after small perturbation of the scores with noise. (vi)~Robust (binned) average rank. (vii)~Elimination rank -- equivalent to an ``Exhaustive Ballot''. All strategies have a high Kendall's correlation with the vanilla mean across tasks ($\tau>0.87$).
Appendix~\ref{app:alt-model-ranking} contains details. Based on these results we choose mean top-1 for VTAB since it is simple, interpretable, can be computed for each task independently, and is highly correlated with more complex appraoches.

\textbf{Representative Subset}\quad
For prototyping, a representative subset of VTAB may be useful. We compute the rank correlation between the mean scores produced by each $\binom{20}{5}$  subsets of five tasks, and the full suite. The top-5 subsets tend to span different domains, but differ to each other.
Appendix~\ref{app:subset-model-ranking} contains the results. Although a subset might be useful for screening models, these were computed on our set of models, and may correlate less well in other experiments. Using subsets also increases the risk of meta-overfitting, which VTAB aims to avoid by having many tasks.

\textbf{Limiting Evaluation Cost}\quad
The cost is near linear in the schedule length. To determine a minimal, yet meaningful, schedule, we sweep over schedules ranging from $40$ to \num{40000} steps, using batch size $512$.
\ref{fig:budget-plot-mean} summarizes the results, details Figs.~\ref{fig:budget-plot-1k} and \ref{fig:budget-plot-full},~\ref{app:budget}.
Most runs reach their optimal performance within \num{1000} steps and do not improve significantly when trained for longer.

Finetuning and evaluating a given ResNet-50 model on a single Nvidia \texttt{P100} GPU (batch size of 64 images, 0.01 initial learning rate) takes 3 hours for VTAB-1k. To reduce the time and cost we use Google Cloud TPU-v3-16 accelerators and verify that we obtain similar resuts in this setting.

\section{Related Work \label{sec:related_work}}

\textbf{Vision Benchmarks}\quad
The Visual Decathlon~\citep{rebuffi2017} contains ten classification tasks: Omniglot, and nine using natural images. Four overlap with VTAB, but VTAB includes several other domains and tasks. Importantly, these benchmarks have opposite evaluation protocols: The Visual Decathlon allows direct joint optimization on the tasks, but forbids external data. By contrast, VTAB forbids multi-task learning on the downstream tasks, but permits transfer from any arbitrary dataset --- this simulates one-off training of representations which are then applied to novel tasks. Further, VTAB considers a low-sample regime (1000-examples), which reflects performance under a reasonable labelling budget. We conduct experiments showing that the methods ranked according to  VTAB are more likely to transfer to new tasks, than those ranked according to the Visual Decathlon (\ref{sec:visual_decathlon}).

The Facebook AI SSL challenge~\citep{fai2019} was proposed to evaluate self-supervision. It contains four tasks on two datasets of natural images, including classification, detection, and low-label classification. The original paper~\citep{goyal2019scaling} also includes navigation and surface normal estimation. VTAB uses classification tasks only to admit task-independent implementations, and attains diversity with many alternative domains and semantics (localization, counting, etc.). Our study includes self-supervision, but are not restricted to it. Indeed, one challenge is to outperform supervised ImageNet on VTAB. Most importantly, Facebook SSL requires transfer via a shallow network stacked on a fixed CNN, whereas VTAB permits any transfer, and focuses on performance with few samples.

Meta-Dataset~\citep{triantafillou2019meta} contains natural image classification tasks. The protocol differs to VTAB: most of Meta-Dataset's evaluation datasets are also in the training set, and the train/test split is made across classes. Meta-Dataset is designed for few-shot learning, rather than 1000 examples, which may entail different solutions.

\textbf{Representation Learning Evaluation}\quad
Popular evaluations for representation learning are linear/MLP and semi-supervised. In linear/MLP evaluation, widely used in for self-supervised representations~\citep{doersch2015unsupervised,zhang2016colorful,noroozi2016unsupervised,doersch2017multi}, the weights of a pre-trained network are frozen, and a linear layer/MLP is trained on top to predict the labels. Evaluation is often performed on the same dataset that was used to train the representations, but using all labels, defeating the goal of sample efficiency. A semi-supervised protocol is more realistic and performs better~\citep{zhai2019s4l}. This evaluation is performed on a single dataset, by contrast VTAB concerns transfer to \emph{unseen} tasks. Linear evaluation is sensitive to small upstream training details, such as the ResNet version or regularization, that make little difference when training end-to-end~\citep{kornblith2018better,kolesnikov2019revisiting}. Other intrinsic scores have been proposed to measure the disentanglement of representations~\citep{locatello2018challenging}, or the mutual information between the input and representations~\citep{hjelm2018learning}. Both were shown to be weakly correlated with representation utility~\citep{locatello2018challenging,tschannen2019mutual}.

\textbf{Generative Evaluation}\quad
Evaluation of generative quality is a difficult goal with numerous proxy metrics, such as reconstruction error, FID~\citep{heusel2017gans}, IS~\citep{salimans2016improved}, precision-recall~\citep{sajjadi2018assessing}, log likelihood (for auto-regressive models), or the ELBO (for VAEs). Sometimes generative models are also evaluated using linear classification~\citep{radford2016,donahue2016adversarial,dumoulin2016adversarially}, or semi-supervised learning~\citep{kingma2014semi, narayanaswamy2017learning,tschannen2018recent}.

\section{Discussion}
Our study answers important and timely questions on representation learning. First, how effective are supervised ImageNet representations? ImageNet labels are indeed effective for natural tasks. Further, although the literature reports mixed results~\citep{liu2017detecting,raghu2019,neumann2019domain}, we find that they also work for specialized tasks despite the domain-shift (medical or remote sensing). However, for structured understanding, supervised representations can be poorer than unsupervised ones. As future work, training on diverse data sources, such as video, may enable moving beyond the ``ImageNet-like'' tasks.

Second, how do representations trained via generative and discriminative models compare? Generation is a useful task in of itself, but generative losses, in their current form, seem less promising as means towards learning how to represent data. BigBiGAN is a notable recent exception, that is on par with self-supervision and warrants more exploration.

Third, to what extent can self-supervision replace labels? Self-supervised learning appears promising, even outperforming supervision on some tasks that require structured understanding. Interestingly, self-supervision can almost (but not quite) replace $90\%$ of ImageNet labels; the gap between pre-training on $10\%$ labels with self-supervison, and $100\%$ labels, is small. Further, self-supervision adds value on top of ImageNet labels \emph{on the same data}. Overall, it seems we are quite far from general visual representations; simply adding more data on the \taskSpecialized{} and \taskStructured{} tasks is better than the pre-training strategies we evaluated. However, self-supervison can be applied to any dataset of images or videos, so combining large scale open-domain self-supervision with ImageNet or other label sources may be promising future work.

VTAB measures representation quality as the ability to adapt with few examples to diverse, unseen tasks. Here, we provide a large study of upstream training losses and control other factors such as hyperparameter sweep, architecture, transfer algorithm, preprocessing, and pre-training data. However, VTAB may be used to analyze and optimze many factors involved in learning generalizable representations. Varying other factors to improve VTAB is valuable future research. To isolate the effect of each factor, confounders such as hyperparameter sweep size should be controlled. The only approach that is out-of-bounds is to condition the algorithm explicitly on the VTAB tasks, since this would compromise representation generalization.

Code and data to run VTAB are made available, and progress is monitored at~\url{github.com/google-research/task_adaptation}.
The models presented here are released. We hope that VTAB drives representation learning towards making deep learning accessible to the many problems without vast labelling budgets.

\subsubsection*{Acknowledgments}
We thank Raphael Marinier and Anton Raichuk for help building the DMLab classification task, and Samy Bengio and Kevin Murphy for useful discussions.

\end{document}
